\section{Introduction}\label{sec:intro}

\subsection{Preliminaries}\label{sec:prelim}

All graphs are finite and simple unless stated otherwise. In a loopless
multigraph, parallel edges are distinct and contribute to degrees with
multiplicity. For a graph $G$, write $V(G)$, $E(G)$ and $\Delta(G)$ for
its vertex set, edge set and maximum degree. For $v\in V(G)$, let
$N_G(v)$ be its set of neighbors and let $d_G(v)$ be its degree.
For $X\subseteq V(G)$, write $G[X]$ for the subgraph induced by $X$ and
$d_G(v,X)$ for the number of edges from $v$ to $X$. Let $N_G[v]=N_G(v)\cup\{v\}$.
%We take $\Delta(G)=0$ for a graph with no vertices. We use $\Delta$ for the maximum degree of a fixed graph and $D$ when retaining the original degree bound after restricting the graph. 
A graph is
$d$-regular if all its vertices have degree $d$. An ordered edge is an
edge together with an ordering of its endpoints.
Write $C_n$ for the cycle on $n\geq 3$ vertices, $K_n$ for the complete
graph on $n$ vertices, and $K_{a,b}$ for the complete bipartite graph with
parts of sizes $a$ and $b$.

A vertex cover of $G$ is a set meeting every edge, and $\tau(G)$ is
the minimum size of a vertex cover. An independent set is a set of
pairwise nonadjacent vertices. For a subgraph $H$ of $G$ with at least
one edge, its Ore-degree in $G$ is
$\max_{uv\in E(H)}(d_G(u)+d_G(v))$~\cite{FaronPostle2019}.
We put $x_+=\max\{x,0\}$ for real $x$.

The \emph{underlying simple graph} of a loopless multigraph $G$, also
called its \emph{simple support}, has vertex set $V(G)$ and one edge
$uv$ for each pair of distinct vertices $u,v$ joined by an edge of $G$.
The line graph $L(G)$ has
one vertex for each edge of $G$, with two distinct edges adjacent
precisely when they share an endpoint in $G$. For a graph $H$,
its square $H^2$ has vertex set $V(H)$, with two distinct vertices
adjacent precisely when $H$ contains a path of length at most two
between them. We write $\omega(H)$ for the maximum size of a clique in
$H$, that is, a set of pairwise adjacent vertices. A set
$S\subseteq E(G)$ is a \emph{strong clique} in $G$ if it is a clique
in $L(G)^2$.
%Equivalently, any two distinct edges of $S$ share an endpoint or have endpoints joined by an edge of $G$.

For a positive integer $k$, write $[k]=\{1,\ldots,k\}$. A proper vertex
$k$-coloring of $H$ is a map $V(H)\to[k]$ assigning distinct colors to
adjacent vertices. The chromatic number $\chi(H)$ is the least $k$ for
which $H$ admits such a coloring, with $\chi(H)=0$ when $V(H)$ is empty.
A proper vertex coloring of $L(G)^2$ gives a \emph{strong edge
coloring} of $G$. The minimum number of colors in such a coloring is
the \emph{strong chromatic index} $\chi(L(G)^2)$.

A real symmetric matrix $M$ is positive semidefinite, written
$M\succeq0$, if $x^{\mathsf T}Mx\geq0$ for every real column vector
$x$ of the appropriate size. The trace $\tr(M)$ of a square matrix
$M$ is the sum of its diagonal entries.

\subsection{Background and Main Result}\label{sec:background}

In 1985, at a seminar in Prague, Erd\H{o}s and Ne\v{s}et\v{r}il asked
how large the strong chromatic index can be for a graph of given
maximum degree~\cite{FGST1989,ErdosNesetril1989}. Erd\H{o}s~\cite{Erdos1988}
stated their conjecture with the bound $\frac{5}{4}\Delta(G)^2$.
Faudree, Schelp, Gy\'arf\'as and Tuza~\cite{FGST1990} formulated a
version for loopless multigraphs that gives a sharper bound for odd
maximum degree. For a nonnegative integer $d$, let
$f(d)=d^2+\lfloor\frac{d}{2}\rfloor^2$. Thus $f(d)=\frac{5d^2}{4}$
when $d$ is even and $f(d)=\frac{5d^2-2d+1}{4}$ when $d$ is odd.

\begin{conjecture}[\cite{FGST1990}]
\label{conj:strong-edge-coloring}
Every finite loopless multigraph $G$ satisfies
$\chi(L(G)^2)\leq f(\Delta(G))$.
\end{conjecture}

Faudree, Gy\'arf\'as, Schelp and Tuza also asked how large a set of edges
forced to have distinct colors can be. Hence, they proposed the following
relaxation of \textbf{Conjecture~\ref{conj:strong-edge-coloring}}.

\begin{conjecture}[\cite{FGST1990}]\label{conj:fgst}
Every finite loopless multigraph $G$ satisfies
$\omega(L(G)^2)\leq f(\Delta(G))$.
\end{conjecture}

In this paper, we prove \textbf{Conjecture~\ref{conj:fgst}}.

\begin{theorem}\label{thm:multigraph}
Every finite loopless multigraph $G$ satisfies
$\omega(L(G)^2)\leq f(\Delta(G))$.
\end{theorem}

Blow-ups of a five-cycle attain the bound. Let $C_5=v_0v_1v_2v_3v_4v_0$. Replace each $v_i$ by an independent set $V_i$ and add every edge between $V_i$ and $V_{i+1}$, with indices taken modulo five. It is easy to see that the entire edge set of the blow-up is a strong clique. Taking $|V_i|=t$ for each $i$ gives
maximum degree $2t$ and $5t^2=f(2t)$ edges, as illustrated for
$t=2$ in Figure~\ref{fig:blowup}. Taking $|V_2|=|V_3|=r+1$ and
$|V_0|=|V_1|=|V_4|=r$ gives maximum degree $2r+1$ and
$5r^2+4r+1=f(2r+1)$ edges, as required.

\begin{figure}[!htbp]
\centering
\begin{tikzpicture}[scale=0.9,every node/.style={font=\small}]
\foreach \i in {0,...,4}
  \coordinate (v\i) at ({90+72*\i}:1.16);
\foreach \i/\j in {0/1,1/2,2/3,3/4,4/0}
  \draw[unmarked edge] (v\i)--(v\j);
\foreach \i in {0,...,4}
  \node[graph vertex] at (v\i) {};
\node at (0,-2.08) {$C_5$};

\begin{scope}[xshift=5.45cm]
\foreach \i in {0,...,4}{
  \pgfmathsetmacro{\angle}{90+72*\i}
  \coordinate (a\i) at (\angle:1.18);
  \coordinate (b\i) at (\angle:1.76);
  \coordinate (part\i) at (\angle:1.47);
  \begin{scope}[shift={(part\i)},rotate=\angle]
    \draw[draw=black!45,fill=black!7] (0,0)
      ellipse[x radius=.42,y radius=.17];
  \end{scope}
}
\foreach \i/\j in {0/1,1/2,2/3,3/4,4/0}{
  \draw[marked edge] (a\i)--(a\j) (a\i)--(b\j)
    (b\i)--(a\j) (b\i)--(b\j);
}
\foreach \i in {0,...,4}{
  \node[graph vertex] at (a\i) {};
  \node[graph vertex] at (b\i) {};
}
\node at (0,-2.08) {$t=2$};
\end{scope}
\end{tikzpicture}
\caption{The cycle $C_5$ and its blow-up with two vertices per part.}
\label{fig:blowup}
\end{figure}

The proof is computer-assisted. We verify the flag inequalities used
in Sections~\ref{sec:uniform} and~\ref{sec:exact} by checking rational
certificates in exact arithmetic. For several small values of $\Delta$,
we derive necessary conditions for a strong clique to have more than
$f(\Delta)$ edges. In some cases, we enumerate the possible configurations
and verify that none can contain such a clique. In the other cases, we
encode the necessary conditions in Boolean formulas and verify that
none of these formulas has a satisfying assignment. The appendix describes
these checks. \textbf{Appendix~\ref{app:reproduction}} gives the commands
for reproducing them with the data, certificates and verification programs
in the repository~\cite{FGSTCode}.

